\documentclass[10pt,a4paper]{amsart}
\usepackage[margin=19mm]{geometry}
\usepackage{amsmath,amssymb,mathtools}
\usepackage[scr=rsfs,cal=euler]{mathalfa}
\usepackage{xfrac}
\usepackage[unicode,hidelinks,bookmarks=false]{hyperref}
\newcommand{\ie}{i.e.\ }
\newcommand{\cf}{cf.\ }
\newcommand{\resp}{respectively\ }
\newcommand{\Subsection}{\subsection}
\providecommand{\nobreakdash}{\nobreak\hbox{-}}
\setlength{\emergencystretch}{2em}
\multlinegap0pt
\usepackage{amssymb}
\usepackage{bm}
\usepackage[shortlabels]{enumitem}
\newtheorem{lemma}{Lemma}
\newtheorem{theorem}{Theorem}
\def\D{\mathbb{D}}
\def\HH{{\mathcal H}}
\def\YY{{\mathcal Y}}
\renewcommand{\leq}{\leqslant}
\title{Littlewood subordination for de Branges--Rovnyak spaces}
\author{Emmanuel Fricain, Andreas Hartmann, William T. Ross, Dan Timotin}
\date{Source: arXiv:2609.10115v1 (CC BY 4.0)}
\begin{document}
\maketitle

\noindent\emph{Source and licence.} Littlewood subordination for de Branges--Rovnyak spaces. Version arXiv:2609.10115v1 (CC BY 4.0), distributed by arXiv under the Creative Commons Attribution 4.0 International licence, http://creativecommons.org/licenses/by/4.0/, as recorded in the arXiv licence metadata for this exact version and shown on the abstract page https://arxiv.org/abs/2609.10115v1. Full licence notice: \texttt{licenses/arxiv-2609.10115-CC-BY-4.0.txt}.\par\smallskip

\noindent\textit{Editorial scope.} Counted unit: the article's theorem inner (source0.tex lines 1069--1079). The article's complete original proof of it is delivered with it (source0.tex lines 1081--1145, one \texttt{proof} environment of 65 lines). No mathematics is written, repaired or re-derived: the statement and the proof are the article's own lines, in the article's own notation, and no retained line is substituted or re-flowed.  Retained uncounted as the source-local context the proof needs: lines 1023--1029; lines 1063--1065.  Nothing was omitted from any retained span, so the package carries no omission record.  No retained line contains a citation, so the wrapper carries no bibliography and no bibliography entry.  No retained line contains a figure environment, an image-inclusion command or drawing code, so no image is delivered and the package declares no asset dependency.  Editorial formatting only, in addition: an AMS \texttt{amsart} preamble that restates the article's own declarations and macros used by the retained lines, namely the environments \texttt{lemma}, \texttt{theorem} on separate counters, together with the standard packages those lines need.  The retained environments print in this document as Lemma 1, Theorem 1 in order of appearance, whereas the article prints the same items with the numbering of its own full document.  The complete native source of the article as distributed in the arXiv e-print for this exact version is bundled unchanged as \texttt{sources/209875/source0.tex}.
\begin{lemma}\label{le:mistery identity}
	Let $\phi$ be inner with $\phi(0)=0$ and $\alpha\in H^\infty$. Then
	\begin{equation}\label{eq:mystery}
			T_{\overline{\alpha\circ\phi}}C_\phi=C_\phi T_{\overline\alpha}
	\end{equation}
Consequently, $\YY_\phi^\perp = H^2 \ominus \mathcal{Y}_{\phi}$ is invariant with respect to multiplication by the function $\alpha\circ\phi$.
\end{lemma}

\begin{equation}\label{eq:norm in H(b)}
\| f \|_{\HH(b)}^2=\sup_{g\in H^2}(\|  f+bg \|^2-\|  g \|^2).
\end{equation}

\begin{theorem}\label{th:compo inner}
	Let $b \in H^{\infty}_{1}$ and $\phi$ be an inner function such that $\phi(0) = 0$. Then we have the following. 
\begin{enumerate}
    \item  $C_\phi$ is an isometry from $\mathcal{H}(b)$ into $\mathcal{H}(b \circ \phi)$.
    \item We have 
    $$C_\phi(\HH(b))=\HH(b\circ\phi)\cap\YY_\phi$$ and the space $\HH(b\circ\phi)$ can be orthogonally decomposed as
		\begin{equation}\label{eq:decomp H(b circ phi)}
			\HH(b\circ\phi)= C_\phi(\HH(b))\oplus_{	\HH(b\circ\phi)} 	\left(\HH(b\circ\phi)\cap\YY_\phi^\perp\right).
		\end{equation}
\end{enumerate}		 
\end{theorem}

\begin{proof}
For an arbitrary $F\in \HH(b\circ\phi)$, decompose it orthogonally in $H^2$ as 
$$F=f\circ \phi+u, \quad f\circ \phi\in \YY_\phi, \quad u\in\YY_\phi^\perp.$$ From \eqref{eq:norm in H(b)} we have
	\[
		\|F\|^2_{\mathcal{H}(b \circ \phi)}=\sup_{H\in H^2}\left( \|F+(b\circ \phi) H\|^2-\|H\|^2\right).
	\]
	Decomposing $H \in H^2$ as 
    $$H=h\circ \phi+v, \quad h\circ \phi\in \YY_\phi, \quad v\in\YY_\phi^\perp,$$ we have, again using \eqref{eq:norm in H(b)},
	\[
	\begin{split}
		\|F\|^2_{\mathcal{H}(b \circ \phi)}&= \sup_{\substack{h \in H^2 \\ v \perp \YY_{\phi}}} \left(\| f\circ\phi+u+(b\circ\phi)(h\circ \phi+v)\|^2-\|h\circ\phi+v\|^2\right)\\
		&=\sup_{\substack{h \in H^2 \\ v \perp \YY_{\phi}}} \left(\| (f+bh)\circ\phi+u+(b\circ\phi)v\|^2-\|h\circ\phi\|^2-\|v\|^2\right).
	\end{split}
	\]
	By Lemma~\ref{le:mistery identity}, $u+(b\circ \phi)v\in \YY_\phi^\perp$, and so the equality above continues as 
\begin{align*}
& = \sup_{\substack{h \in H^2 \\ v \perp \YY_{\phi}}} \left(\| (f+bh)\circ\phi\|^2+\|u+(b\circ\phi)v\|^2-\|h\circ\phi\|^2-\|v\|^2\right)\\
	&=\sup_{\substack{h \in H^2 \\ v \perp \YY_{\phi}}} \left(\left(\| (f+bh)\circ\phi\|^2-\|h\circ\phi\|^2\right) + \left(\| u+(b\circ \phi)v\|^2-\|v\|^2\right)\right).
\end{align*}
	The first parentheses in the previous line depend only on $ h $, while the second depends only on $ v $. Thus,  we actually have
		\begin{align}\label{eq:suprema}
			\|F\|^2_{\mathcal{H}(b \circ \phi)} & =
		\sup_{h\in H^2} \left(\| (f+bh)\circ\phi\|^2-\|h\circ\phi\|^2\right)\\
		& \quad \quad  +\sup_{v\perp\YY_\phi} \left(\| u+(b\circ \phi)v\|^2-\|v\|^2\right).\nonumber
		\end{align}
Note that both suprema on the right hand side of the previous line are bounded above by $\|F\|^2_{\mathcal{H}(b \circ \phi)}$.


	Since $C_\phi$ acts isometrically on $H^2$, the first supremum in~\eqref{eq:suprema} is
	\begin{equation}\label{eq:first supremum}
			\sup_{h\in H^2} \left(\| (f+bh) \|^2-\|h \|^2_2\right)=\|f\|^2_{\HH(b)}.
	\end{equation}
	In particular, we see that if  $F =   f\circ\phi \in \HH(b\circ\phi) $, then $f \in \HH(b)$ ($u = 0$) and
	\[
	\|f\|^2_{\HH(b)}=\|C_\phi f\|^2_{\HH(b\circ\phi)},
	\]
	which proves (a).
	
	
 For the second supremum on the right hand side of \eqref{eq:suprema}, we know that 
	\[
	\| u\|_{\HH(b\circ\phi)}^2=\sup_{h \in H^2}\big( \| u+(b\circ\phi)h\|^2-\|h\|^2\big).
	\]
Now write $h=g\circ \phi+w$ and again use the invariance of $\YY_\phi^\perp$ with respect to multiplication by $b\circ\phi$ (Lemma~\ref{le:mistery identity})  to obtain the orthogonal decompositions
	\[
	\begin{split}
			\| u\|_{\HH(b\circ\phi)}^2&=\sup_{g,w} \left(\| (bg)\circ\phi+u+(b\circ \phi)w\|^2-\|w+g\circ\phi\|^2\right)\\
&		=\sup_{g,w} \left(\|(bg)\circ\phi\|^2+\| u+(b\circ\phi)w\|^2-\|w\|^2-\|g\circ\phi\|^2\right)\\
&=\sup_{g,w}\left((\|bg\|^2-\|g\|^2)+\left(\| u+(b\circ\phi)w\|^2-\|w\|^2\right)\right). 
	\end{split}
		\]
	Since $|b| \leq 1$ on $\D$, the expression in the first parentheses of the last equation is nonpositive, and so we obtain
		\[
			\| u\|_{\HH(b\circ\phi)}^2\leq \sup_w \left(\| u+(b\circ\phi)w\|^2-\|w\|^2\right) .
		\]
		In fact, we have 
	\begin{equation}\label{eq:second supremum}
			\| u\|_{\HH(b\circ\phi)}^2=\sup_w \left(\| u+(b\circ\phi)w\|^2-\|w\|^2\right)
	\end{equation}
	since we may choose $g=0$ in the above. Thus,  $ u\in \HH(b\circ\phi) $. Moreover, combining \eqref{eq:suprema}, \eqref{eq:first supremum}, and \eqref{eq:second supremum}, we obtain the equality
	\[
		\|F\|^2_{\mathcal{H}(b \circ \phi)}=\| f\circ\phi\|^2_{\HH(b\circ\phi)}+	\| u\|_{\HH(b\circ\phi)}^2.
	\]
	It follows that the orthogonal decomposition $F=f\circ\phi+u$  in $H^2$  is also an orthogonal decomposition in $\HH(b\circ\phi)$. This verifies \eqref{eq:decomp H(b circ phi)} and thus completes the proof. 
	\end{proof}

\end{document}
